\section{Conclusion}
\label{sec:conclusion}

This paper evaluated YOLO26m, an NMS-free single-stage detector, on a
publicly available three-class vehicle detection dataset (9\,725 images;
\textit{bus}, \textit{car}, \textit{van}) exported from Roboflow Universe
under CC~BY~4.0 \cite{roboflow_dataset}.
We reported detection performance, latency, and a TIDE-style error
taxonomy for the deployed checkpoint, documented a silent label-format
corruption that affects training data assembled from mixed annotation
sources, and demonstrated an end-to-end deployment pipeline on consumer
Apple Silicon hardware.

% -------------------------------------------------------------------------
\subsection*{Summary of Findings}
% -------------------------------------------------------------------------

\begin{enumerate}

  \item \textbf{YOLO26m achieves competitive accuracy on a heterogeneous
    traffic dataset.}
    The deployed checkpoint reaches mAP$_{50}$~=~0.784 and
    mAP$_{50:95}$~=~0.594 on the held-out test split, with precision
    0.845 and recall 0.752.
    Per-class analysis shows that \textit{bus} and \textit{car} are
    detected reliably (AP$_{50}$ 0.800 and 0.786, respectively), while
    \textit{van} is harder due to visual overlap with \textit{car}
    (AP$_{50}$ 0.766, dominant confusion pair: van$\rightarrow$car,
    14 cases).

  \item \textbf{Missed detections are the primary error mode.}
    Of the 413 ground-truth instances in the test set, 78 (19\,\%) were
    missed at a confidence threshold of 0.40.
    Background false positives (22) and classification errors (23) are
    secondary concerns.
    The model's precision-recall profile suggests that the confidence
    threshold can be tuned downward to improve recall at the cost of
    additional false positives — a trade-off that depends on the
    application.

  \item \textbf{NMS-free inference eliminates post-processing latency.}
    Post-processing contributes only 0.08\,ms per image; the total
    end-to-end CPU latency is 99.1\,ms (10.1\,FPS), dominated by the
    96.31\,ms forward pass.
    At this throughput the system is suitable for frame-sampled analysis
    or offline batch processing; MPS acceleration on the Apple M4 Pro is
    expected to bring throughput above 20\,FPS.

  \item \textbf{Mixed bounding-box/polygon label files are a silent
    source of training corruption.}
    Three of the 9\,218 training label files contained mixed-format rows.
    Ultralytics~8.4.53 silently reinterprets these files entirely as
    polygon annotations, producing incorrect bounding boxes without any
    warning.
    A lightweight token-count audit (released in
    \texttt{trace\_experiments.py}) detects and corrects the problem in
    seconds.

  \item \textbf{The deployment pipeline integrates detection, logging, and
    alerting on consumer hardware.}
    The VehicleScope~AI Streamlit application, n8n webhook, Google Sheets
    logger, and Telegram alerter form a complete, low-cost deployment
    stack that runs without cloud GPU infrastructure.

\end{enumerate}

% -------------------------------------------------------------------------
\subsection*{Future Work}
% -------------------------------------------------------------------------

The findings reported here are based on a single pre-trained checkpoint
and exclude the planned re-training experiments.
The following directions will complete and extend the study.

\begin{description}

  \item[Baseline comparison and augmentation ablation (Tables T8, T9).]
    Training YOLO26m, YOLOv8m, YOLO11m, and RT-DETR-l from COCO weights
    on the cleaned dataset, over three seeds, will provide a statistically
    grounded comparison.
    The augmentation ablation (photometric only vs.\ Ultralytics default
    vs.\ photometric~+ light geometric) will quantify the sensitivity of
    YOLO26m to data augmentation choices.

  \item[Quantifying the label-fix effect.]
    Re-training with the original (uncleaned) labels and comparing against
    the cleaned-label model will isolate the accuracy impact of the
    mixed-format bug.
    The magnitude of the effect depends on how heavily the corrupted files
    are sampled during training.

  \item[Cross-domain evaluation (Surabaya dataset).]
    Collecting and annotating a dataset of Surabaya, Indonesia, traffic
    footage — split by recording location to prevent data leakage — will
    allow measurement of the Philippine-to-Indonesian transfer gap and
    evaluation of joint training.
    This dataset is planned in \texttt{trace\_novelty.py} and will be
    released under an open licence.

  \item[MPS latency benchmark.]
    Running the latency stage with \texttt{--device mps} on the Apple M4
    Pro will provide a complete picture of the deployment cost and
    enable comparison with the CPU baseline.

  \item[Confidence-threshold sensitivity.]
    A sweep of the confidence threshold from 0.20 to 0.60 will
    characterise the precision-recall trade-off and identify the
    operating point that minimises under-counting for traffic-management
    use cases.

  \item[Prevalence audit across Roboflow Universe.]
    Applying the \texttt{stage\_audit} script to a random sample of
    public vehicle datasets on Roboflow Universe will estimate how
    widespread the mixed-format problem is in community datasets.

\end{description}

% -------------------------------------------------------------------------
\subsection*{Data and Code Availability}
% -------------------------------------------------------------------------

The experiment harness (\texttt{trace\_experiments.py},
\texttt{trace\_novelty.py}), the CI workflow
(\texttt{.github/workflows/ci.yml}), the test suite (\texttt{tests/}),
and the Streamlit application (\texttt{app/}) are available in the public
GitHub repository \texttt{indri007/TRACE}.
The dataset is available from Roboflow Universe at
\url{https://universe.roboflow.com/personal-project-kej16/vehicle-detection-vznzd-dkl8g/dataset/1}
under CC~BY~4.0 \cite{roboflow_dataset}.
Model weights and large data files are excluded from the repository;
instructions for reproducing all results are provided in
\texttt{README.md}.
